%% preface.tex --- preface and reading guide

\chapter*{Preface}
\addcontentsline{toc}{chapter}{Preface}
\markboth{Preface}{Preface}

Machine learning is taught twice over. There is the version made of intuition and
pictures, which is pleasant to listen to and impossible to reproduce, and there is
the version made of derivations, which is reproducible and, presented badly,
unreadable. This book tries to be one thing: the intuition, and then the algebra
that makes the intuition precise, in that order, with nothing skipped in between.

The rule I set myself while writing was that the reader should never meet the
phrase \emph{it can be shown that}. Where a result is used, it is derived. Where a
number is quoted --- the value of a cost function, the output of three iterations
of an algorithm, an error rate on held-out data --- it was computed by a script
that ships with the book, not estimated for the sake of a clean-looking page. If
you work an example by hand and get a different answer, one of us is wrong, and
the printed digits are precise enough to tell which.

\section*{Who this book is for}

Undergraduates meeting machine learning for the first time, and self-learners with
the same preparation. I assume you can program in Python, that you have seen data
structures and algorithms, and that you have had one course in probability and
statistics. I do \emph{not} assume linear algebra or multivariable calculus.
Appendix~\ref{ch:appA} builds the small amount of both that the book needs, and
each chapter points into it at the moment the tool is first required, so you can
learn the mathematics on demand rather than in a block at the start.

\section*{How the book is organized}

Parts I and II develop the whole machinery of supervised learning on the simplest
possible models: linear regression, then logistic regression. Almost every idea in
the rest of the book --- cost functions, gradient descent, overfitting,
regularization, evaluation --- appears there first, where it can be drawn on a
page and checked with a calculator. Part III adds neural networks, which is
mostly a matter of stacking what Part II already established. Part IV is about
making a learning system work: diagnosing it, improving it, and measuring it
honestly. Parts V through VII cover unsupervised learning, recommender systems,
and reinforcement learning. Part VIII treats attention, transformers, and large
language models. Part IX is about everything that happens after the model is
accurate: designing, building, shipping, and operating a system that contains one.

Three reading paths through this material:

\begin{description}
  \item[A first course.] Chapters 1--17 in order, one chapter per lecture-sized
    session, with the labs. Chapters 12, 18 and Part IX become reading and a
    term project.
  \item[Revision.] Each chapter's Learning Objectives, Key Insight and Pitfall
    boxes, the chapter summary, then the exercises. If an exercise is hard, the
    section it came from is the section to reread.
  \item[Practitioner.] Chapters 1--7 and 10--11 for the vocabulary, then
    Part IX, then back to whichever of Parts III--VIII your problem needs.
\end{description}

\section*{How to read a chapter}

Every chapter opens with Learning Objectives and closes with a summary and five
kinds of exercise: short questions, analytical problems that ask you to derive or
compute, design problems with no single right answer, programming exercises, and
multiple-choice items for self-testing. Boxes carry a fixed meaning throughout:

\begin{description}
  \item[Definition] the precise statement of a term, worth memorising.
  \item[Key Insight] the one idea to carry out of a section.
  \item[Derivation] algebra, step by step; read it with a pen.
  \item[Worked Example] numbers, carried through by hand.
  \item[Algorithm] pseudocode you could implement directly.
  \item[Pitfall] a mistake that is common, quiet, and examinable.
  \item[Research Spotlight] the paper an idea came from, in one page.
  \item[In the Lab] the point at which to stop reading and write code.
  \item[Ethics and Impact] the consequences of deploying what you just built.
\end{description}

The labs are not optional decoration. Machine learning is a subject in which it is
entirely possible to follow every derivation and still be unable to fit a model,
because the difficulties in practice --- a transposed array, a leaked test set, a
step size off by a factor of ten --- are invisible in the algebra. Each lab
implements a method twice, first as an explicit loop that matches the equations
line for line, then in vectorized form, and checks that the two agree.

\section*{Acknowledgements}

This book grew out of lecture notes for a one-semester undergraduate course, and
it owes an evident debt to the way Andrew Ng has taught this material to a very
large number of people. Its statistical framing follows James, Witten, Hastie and
Tibshirani; where it goes deeper it usually follows Bishop. The errors that remain
are mine, and I would be glad to hear about them.

\vspace{1cm}
\noindent\textit{Rokan Uddin Faruqui}\\
\noindent\textit{2026}

\cleardoublepage
